\documentclass[11pt]{article}
\usepackage[T1]{fontenc}
\usepackage{lmodern}
\usepackage[margin=1in]{geometry}
\usepackage{amsmath,amssymb,amsthm,mathtools}
\usepackage[colorlinks=true,linkcolor=blue,citecolor=blue,urlcolor=blue]{hyperref}
\usepackage{microtype}
\newtheorem{theorem}{Theorem}[section]
\newtheorem{lemma}[theorem]{Lemma}
\newtheorem{remark}[theorem]{Remark}
\newcommand{\T}{\mathbb T}
\newcommand{\R}{\mathbb R}
\newcommand{\PP}{\mathbb P}
\newcommand{\diver}{\operatorname{div}}
\newcommand{\esssup}{\operatorname*{ess\,sup}}
\newcommand{\osc}{\operatorname*{ess\,osc}}
\newcommand{\dd}{\,\mathrm d}
\title{A deterministic H\"older exponent for parabolic SPDEs\protect\\with smooth transport noise}
\author{Kaiwen Jin\thanks{Contact: \texttt{yui.ui.siki@gmail.com}.}}
\date{October 15, 2026}

\begin{document}
\maketitle
\begin{abstract}
We consider scalar divergence-form parabolic It\^o equations on the torus,
with bounded measurable diffusion and progressively measurable transport
coefficients that are smooth in space with deterministic derivative bounds.
Under stochastic parabolicity, we prove that every variational solution with
deterministic smooth initial data has a positive deterministic parabolic
H\"older exponent on each compact positive-time strip. The H\"older norm may
remain random and unbounded. The stochastic-flow transformation of Agresti,
Sauerbrey and Veraar removes the noise, but its global ellipticity ratio can
depend on the realization. We normalize the continuous inverse flow derivative
by a constant affine map near each spacetime point. This gives uniform local
energy-class constants on possibly random neighborhoods. Compactness preserves
the common exponent, and composition with the inverse flow gives the assertion
in the original coordinates.
\end{abstract}

\section{Introduction and main result}

For a scalar deterministic divergence-form parabolic equation, bounded
measurable coefficients and a fixed ellipticity ratio give a common positive
H\"older exponent. Smooth transport noise permits a stochastic change of
coordinates that removes the noise, but the resulting global ellipticity ratio
need not have a deterministic bound. Agresti, Sauerbrey and Veraar
\cite[Theorem 1.1 and Section 1.2]{ASV} obtain pathwise H\"older regularity and
ask whether its exponent can be chosen independently of the realization.
Their Proposition 4.2 exhibits unbounded transformed ellipticity ratios;
Proposition 4.3 gives uniform exponents under additional structure. The
distortion of the transformed coefficients leaves open the regularity of the
original solution.

Here we answer the deterministic-exponent question for the full smooth-noise
class in AIM Problem 455. A large but finite coordinate distortion changes
the H\"older constant, whereas a fixed affine change of coordinates preserves
the exponent. At each point we use the inverse flow derivative itself for this
affine change. Its continuity gives a neighborhood on which the normalized
diffusion has deterministic bounds. The neighborhood can shrink with the
sample, which is harmless when only the exponent must be deterministic.
The stochastic flow and distributional It\^o--Wentzell formula are taken from
\cite{ASV}. The deterministic regularity input is the parabolic De Giorgi
class theorem of Liao \cite{Liao}; we verify its energy inequalities for the
nonsymmetric coefficients and bounded drift that occur here.

Write $\T^d=\R^d/\mathbb Z^d$, with its induced spatial distance
$\operatorname{dist}$. All function spaces below are real. Matrix norms are
operator norms; other fixed matrix norms change only dimension constants.
The space $H^{-1}(\T^d)$ is the dual of $H^1(\T^d)$ with respect
to the $L^2$ pairing. On a filtered probability space, let $(W^k)_{k\ge1}$
be independent standard Brownian motions. Consider the scalar It\^o equation
\begin{equation}\label{eq:spde}
 \dd u=\diver(a\nabla u)\dd t
       +\sum_{k\ge1}(b_k\cdot\nabla u)\dd W^k_t,
 \qquad u(0)=u_0.
\end{equation}
The matrix $a$ and the vector fields $b=(b_k)_{k\ge1}$ are progressively
measurable in $(t,\omega)$ and measurable in space. We assume deterministic
constants $M<\infty$ and $\nu>0$ such that, almost everywhere,
\begin{align}
 |a|+\Big(\sum_k|b_k|^2\Big)^{1/2}&\le M,\label{eq:bound}\\
 \eta^\top\Big(a-\tfrac12\sum_k b_k\otimes b_k\Big)\eta
       &\ge\nu|\eta|^2\qquad(\eta\in\R^d).\label{eq:parabolicity}
\end{align}
For every integer $j\ge1$, $b$ is spatially $C^j$ as an
$\ell^2(\R^d)$-valued map and
\begin{equation}\label{eq:smooth}
 \sup_{t,\omega}\|b(t,\omega,\cdot)\|_{C^j(\T^d;\ell^2(\R^d))}
 \le M_j<\infty
\end{equation}
for fixed deterministic $M_j$. No symmetry of $a$ is required.

We use the variational solution
$u\in C([0,T];L^2)\cap L^2(0,T;H^1)$, almost surely. Thus, for every
smooth spatial test $\varphi$, its $L^2$ pairing satisfies
\[
 (u(t),\varphi)=(u_0,\varphi)
 -\int_0^t(a\nabla u,\nabla\varphi)\dd s
 +\sum_k\int_0^t(b_k\cdot\nabla u,\varphi)\dd W^k_s.
\]
The stochastic integrals are It\^o integrals. A representative below agrees
with this $L^2$-valued process at every time. For $0<\theta<1$,
$C^{\theta/2,\theta}([\tau,T]\times\T^d)$ denotes the bounded continuous
functions whose increments are bounded by a constant times
$|t-s|^{\theta/2}+\operatorname{dist}(x,y)^\theta$.

\begin{theorem}\label{thm:main}
Let $d\ge2$, $T>0$, and let \eqref{eq:bound}--\eqref{eq:smooth} hold.
For every deterministic $u_0\in C^\infty(\T^d)$, the variational solution of
\eqref{eq:spde} has, for every $0<\tau<T$, a representative satisfying
\[
 \PP\bigl(u\in C^{\gamma/2,\gamma}([\tau,T]\times\T^d)\bigr)=1
\]
for a deterministic $\gamma\in(0,1)$ depending only on $d,\nu,M$.
In particular it has the dependence allowed in AIM Problem 455.
The H\"older norm need not be deterministically bounded. The vector fields
may be noncommuting, time-random, and countably infinite in number.
\end{theorem}

The bounds $M_j$ ensure the required regularity of the stochastic flow.
They affect the random neighborhoods and constants in the proof, but not our
choice of exponent. We do not assert moment bounds for the H\"older norm.
We first prove the deterministic normalization lemma, then apply the
stochastic change of coordinates and return to \eqref{eq:spde}.

\section{A local deterministic normalization}

\subsection{Bounded drift at unit scale}
For $R>0$ write $K_R=(-R,R)^d$ and $Q_R=(-R^2,0)\times K_R$.

\begin{lemma}\label{lem:unit}
Suppose $w\in C_tL^2_{\mathrm{loc}}\cap L^2_tH^1_{\mathrm{loc}}$ is a
scalar weak solution on $Q_2$ of
\begin{equation}\label{eq:unit}
 w_t=\diver(H\nabla w)+f\cdot\nabla w.
\end{equation}
Assume $|H|\le\Lambda$, $\xi^\top H\xi\ge\lambda|\xi|^2$ and
$|f|\le1$ almost everywhere, where $\lambda>0$ and $\Lambda<\infty$.
There is $\alpha=\alpha(d,\lambda,\Lambda)\in(0,1)$ such that $w$ has a
$C^{\alpha/2,\alpha}$ representative on $Q_1$. If $w$ has a continuous
$L^2$ trace at time zero and its energy bounds hold to that time, this
representative extends with the same exponent to the upper face of $Q_1$.
Neither $H$ nor $f$ needs to be differentiable, and $H$ need not be symmetric.
\end{lemma}

\begin{proof}
We check the energy-class hypotheses of \cite[Section 1.1]{Liao}.
Fix a level $k\in\R$ and either truncation $z=(w-k)_+$ or $z=(w-k)_-$.
For a nonnegative spatial cutoff $\zeta$, use the signed test
$z\zeta^2$ or $-z\zeta^2$, respectively. Coercivity and the full matrix
norm bound give
\begin{align*}
 \frac12\frac{\dd}{\dd t}\int z^2\zeta^2
 +\lambda\int|\nabla z|^2\zeta^2
 &\le2\Lambda\int |\nabla z|\zeta|\nabla\zeta|z
       +\int |f||\nabla z|z\zeta^2\\
 &\le\frac\lambda2\int|\nabla z|^2\zeta^2
       +\frac{4\Lambda^2}{\lambda}\int z^2|\nabla\zeta|^2
       +\frac1\lambda\int z^2\zeta^2.
\end{align*}
For a time-dependent cutoff, add $\int z^2|\zeta\zeta_t|$ on the right.
All spatial integrals here are over the support of $\zeta$.
These tests are legitimate at the stated weak regularity: apply time
Steklov averaging and smooth Lipschitz approximations to the signed
truncations, integrate on a compact time interval, and pass to the limits.
The averaged gradients and fluxes converge in $L^2$, their products in
$L^1$, and the truncations converge in $H^1$ by the Sobolev chain rule.
Time continuity gives the endpoint $L^2$ terms. No derivative falls on
$H$ or $f$.

For a translated cube $K_R(x)$, choose $\zeta=1$ on $K_\rho(x)$,
supported in $K_R(x)$, with $|\nabla\zeta|\le C_d/(R-\rho)$.
All cubes lie in $K_2$, so $R-\rho\le2$ and the last volume term is
absorbed into a constant times $(R-\rho)^{-2}\int z^2$.
With a time cutoff between $s-h$ and $s-\vartheta$, we obtain
\begin{align}\label{eq:DG1}
 &\esssup_{s-\vartheta<t<s}\int_{K_\rho(x)}z^2
  +\int_{s-\vartheta}^s\int_{K_\rho(x)}|\nabla z|^2\dd t\nonumber\\
 &\qquad\le\Gamma\bigl((R-\rho)^{-2}+(h-\vartheta)^{-1}\bigr)
       \int_{s-h}^s\int_{K_R(x)}z^2\dd t,
\end{align}
where $0<\vartheta<h$ and $\Gamma=\Gamma(d,\lambda,\Lambda)$.
Without a time cutoff, dropping the nonnegative gradient term gives
\begin{equation}\label{eq:DG2}
 \esssup_{s-h<t<s}\int_{K_\rho(x)}z^2
 \le\int_{K_R(x)}z(s-h)^2
  +\frac{\Gamma}{(R-\rho)^2}
       \int_{s-h}^s\int_{K_R(x)}z^2\dd t.
\end{equation}
The initial integral has coefficient exactly one. These inequalities hold
for every admissible nested cylinder, every level, and both signs. Together
with local time continuity, they put $w$ in Liao's class
$\mathfrak B_2(Q_2,\Gamma)$, as defined in
\cite[equations (1.1)--(1.2)]{Liao}.

By \cite[Theorems 2.1 and 5.1]{Liao}, $w$ is locally bounded and
\[
 \osc_{Q_\rho(z_0)}w
 \le C\,\osc_{Q_R(z_0)}w\,(\rho/R)^\alpha,
 \qquad 0<\rho<R,
\]
with $C$ and $\alpha\in(0,1)$ depending only on $d,2,\Gamma$.
Here $Q_R(z_0)$ is the backward cylinder of radius $R$ with top center
$z_0$. Fixed small outer cylinders about points of $Q_1$ have a spatial
and lower-time buffer in $Q_2$. The local boundedness estimate gives a
common finite oscillation bound on these cylinders. For two nearby points,
use a backward cylinder at the later time with radius larger than their
spatial separation and the square root of their time separation.
Oscillation decay gives the parabolic $\alpha$ estimate; boundedness
handles separated points.

For the upper face, first take cylinders with top time $-\varepsilon$.
The outer buffers and the energy bound can be kept fixed as
$\varepsilon\downarrow0$, so the H\"older constants remain bounded.
The representative extends continuously to time zero. Its $L^2$ limit
equals the prescribed trace, by strong $L^2$ continuity. This proves the
last assertion using only backward cylinders.
\end{proof}

\subsection{A continuous matrix field}

\begin{lemma}\label{lem:frame}
Let $K:[0,T]\times\T^d\to\mathrm{GL}(d,\R)$ be continuous, and let
$S_0,c$ be bounded measurable coefficients such that
\[
 |S_0|\le\Lambda_0,\qquad
 \xi^\top S_0\xi\ge\nu|\xi|^2.
\]
Suppose $v\in C([0,T];L^2)\cap L^2(0,T;H^1)$ solves
\begin{equation}\label{eq:frame}
 v_t=\diver(A\nabla v)+c\cdot\nabla v,
 \qquad A=KS_0K^\top.
\end{equation}
Then, for every $0<\tau<T$, $v$ has a representative in
$C^{\alpha/2,\alpha}([\tau,T]\times\T^d)$ for one
$\alpha=\alpha(d,\nu,\Lambda_0)\in(0,1)$.
Bounds on $K,K^{-1}$ and $c$ affect the H\"older constant, but not
this exponent.
\end{lemma}

\begin{proof}
Fix $z_0=(t_0,x_0)$ with $t_0>0$ and set $L=K(z_0)$.
This is a constant invertible matrix. On a sufficiently small neighborhood
$U$ of $z_0$, continuity gives
\[
 P=L^{-1}K,\qquad |P-I|\le\tfrac12.
\]
Consequently
\begin{equation}\label{eq:normalized}
 \xi^\top PS_0P^\top\xi\ge\tfrac\nu4|\xi|^2,
 \qquad |PS_0P^\top|\le\tfrac94\Lambda_0\le4\Lambda_0.
\end{equation}
Choose $r>0$ so small that the affine cylinder
\[
 t=t_b+r^2s,\qquad x=x_0+rLy,\qquad (s,y)\in Q_2,
\]
lies in $U$, lies in a torus coordinate chart, has lower time strictly
positive, and satisfies $r|L^{-1}|\|c\|_\infty\le1$.
If $t_0<T$, take $t_b>t_0$ with
$0<t_b-t_0<\min\{r^2/2,(T-t_0)/2\}$; if $t_0=T$, take $t_b=T$.
Shrinking $r$ makes all these requirements compatible.

For $W(s,y)=v(t_b+r^2s,x_0+rLy)$, the weak equation becomes
\begin{equation}\label{eq:affine}
 W_s=\diver_y(H\nabla_yW)+f\cdot\nabla_yW,
 \qquad H=L^{-1}AL^{-\top}=PS_0P^\top,
 \quad f=rL^{-1}c,
\end{equation}
where coefficients on the right are evaluated at $(t_b+r^2s,x_0+rLy)$.
The constant spatial Jacobian cancels in the weak formulation; the time
factor $r^2$ cancels the two spatial derivative factors and leaves one
factor $r$ on the drift. By \eqref{eq:normalized}, Lemma~\ref{lem:unit}
applies with $\lambda=\nu/4$ and $\Lambda=4\Lambda_0$.
These parameters do not involve $L,r,c$ or the point $z_0$.
Returning by the fixed affine map changes the H\"older constant and
preserves the exponent. At $t_0=T$ we use the upper-trace part of that lemma.

We have obtained the same exponent on a relative neighborhood of every
point of $[\tau,T]\times\T^d$. Local continuous representatives agree on
overlaps, since they agree almost everywhere with the energy solution.
Time continuity identifies each time slice in $L^2$, including the upper
trace, and spatial continuity then identifies it everywhere.
A finite subcover has a positive Lebesgue number $\ell$ for the parabolic
metric. Pairs at distance below $\ell$ are controlled in a common chart
by the maximum of finitely many H\"older constants. For separated pairs,
use $2\|v\|_\infty\ell^{-\alpha}$. Thus the global norm is finite.
Only this norm, the charts and $\ell$ depend on the particular coefficients.
\end{proof}

\section{Stochastic change of coordinates and proof of the theorem}

\subsection{The flow and the transformed equation}

Put
\begin{equation}\label{eq:BS}
 B=\sum_kb_k\otimes b_k,\qquad S=a-\tfrac12B,
 \qquad\Lambda_0=M+\tfrac12M^2.
\end{equation}
Then $|S|\le\Lambda_0$ and $\eta^\top S\eta\ge\nu|\eta|^2$.
Define the vector fields
\begin{equation}\label{eq:muq}
 \mu^i=\tfrac12\sum_{k,j}b_k^j\partial_jb_k^i,
 \qquad q^i=\tfrac12\sum_k(\diver b_k)b_k^i.
\end{equation}
The product rule and Cauchy--Schwarz in $k$ bound their derivatives of
every finite spatial order by finite combinations of $M$ and the $M_j$.
The sums are continuous bilinear contractions of spatially smooth
$\ell^2$-valued maps, and are progressively measurable.

Lift the periodic coefficients to $\R^d$ and solve the It\^o flow equation
\begin{equation}\label{eq:flow}
 \dd\Phi_t(x)=\mu(t,\Phi_t(x))\dd t
             -\sum_k b_k(t,\Phi_t(x))\dd W^k_t,
 \qquad\Phi_0(x)=x.
\end{equation}
We use this It\^o equation directly, so no time regularity or
Stratonovich conversion of the random coefficient $b$ is needed.
The periodic lifts satisfy global spatial Lipschitz and growth bounds.
Their higher spatial derivative bounds verify Assumption 2.5 of
\cite{ASV} with $k=2$ and spatial H\"older parameter one: bounded
third derivatives make second derivatives Lipschitz. For $\mu$,
derivatives of $b$ through order four suffice.
Theorems 2.3 and 2.6 of \cite{ASV} give a $C^2$ diffeomorphism flow
with jointly continuous spatial derivatives. Its inverse
$\Psi_t=\Phi_t^{-1}$ is spatially Lipschitz and temporally
$1/4$-H\"older on compact sets by \cite[Proposition 2.8]{ASV}.
Remarks 2.4 and 2.7 there explicitly allow countably many Brownian motions
and random coefficients merely measurable in time.

These assertions hold on one probability-one event. Periodicity and
pathwise uniqueness give $\Phi_t(x+m)=\Phi_t(x)+m$ for integer $m$.
One may first impose these identities for rational $x$ and integer $m$
at all times, then extend by continuity. The flow therefore descends to
the torus. Define
\begin{equation}\label{eq:FKJ}
 F=D\Phi,\qquad K=F^{-1},\qquad J=\det F.
\end{equation}
The determinant is continuous and nonzero, and $J_0=1$, so $J>0$.
Compactness of $[0,T]\times\T^d$ gives finite pathwise bounds for
$F,K,J,J^{-1}$ and $\nabla\log J$. In particular $K$ is a continuous
invertible matrix field. All these bounds may depend on the realization.

Set $v(t,x)=u(t,\Phi_t(x))$ in the $L^2$ sense. The spatial chain rule
and change of variables give
\begin{align*}
 \|v(t)\|_2^2
 &\le\|\det D\Psi_t\|_\infty\|u(t)\|_2^2,\\
 \|\nabla v(t)\|_2^2
 &\le\|D\Phi_t\|_\infty^2\|\det D\Psi_t\|_\infty
       \|\nabla u(t)\|_2^2.
\end{align*}
Thus $v\in L^2_tH^1$. It is also continuous in $L^2$. Indeed, split
$u(t)\circ\Phi_t-u(s)\circ\Phi_s$ into the pullback of $u(t)-u(s)$
and the change of pullback of the fixed function $u(s)$. The first
term tends to zero by the uniform pathwise pullback bound. For the
second, approximate $u(s)$ in $L^2$ by smooth functions and use uniform
continuity of $\Phi$ and the same bound. This includes both time endpoints.

We apply \cite[Proposition 2.14]{ASV}, the distributional It\^o--Wentzell
formula. Its drift hypothesis is satisfied because
$\diver(a\nabla u)\in L^2_tH^{-1}_{\mathrm{loc}}\subset
L^1_tH^{-1}_{\mathrm{loc}}$, and its noise hypothesis follows from
\[
 \sum_k|b_k\cdot\nabla u|^2\le M^2|\nabla u|^2.
\]
The torus identity lifts to $\R^d$ by periodizing compactly supported
test functions. The formula holds almost surely as an identity of
continuous $H^{-1}_{\mathrm{loc}}$-valued processes, at all times.
With flow coefficients $\mu,-b$, its stochastic integrand is
$(b_k\cdot\nabla u)\circ\Phi-((b_k\cdot\nabla u)\circ\Phi)=0$.
Before spatial pullback, its remaining drift is
\begin{equation}\label{eq:cancellation}
 \diver(a\nabla u)+\mu\cdot\nabla u
 -\sum_k b_k\cdot\nabla(b_k\cdot\nabla u)
 +\tfrac12B:D^2u
 =\diver(S\nabla u)+q\cdot\nabla u.
\end{equation}
For the signs and factors, put $C=2\mu$. Distributionally,
\[
 \sum_kb_k\cdot\nabla(b_k\cdot\nabla u)=B:D^2u+C\cdot\nabla u,
 \qquad
 \diver(B\nabla u)=B:D^2u+(C+2q)\cdot\nabla u.
\]
Substitution proves \eqref{eq:cancellation}. Every differentiated
coefficient in this calculation is a spatially smooth contraction of
$b$; no derivative of $a$ or pointwise Hessian of $u$ is used.

For an $L^2$ vector flux $h$, the distributional change of variables is
\begin{equation}\label{eq:piola}
 (\diver h)\circ\Phi
   =J^{-1}\diver\bigl(JK(h\circ\Phi)\bigr).
\end{equation}
To see the Jacobian factor, the distribution on the left pairs a test
$\varphi$ with $\diver h$ against
$(\varphi/J)\circ\Psi$. At $y=\Phi(x)$ its gradient is
$K(x)^\top\nabla_x(\varphi/J)(x)$. Changing variables proves
\eqref{eq:piola}, first for smooth $h$ and then for $L^2$ fluxes by
density. Since $(\nabla u)\circ\Phi=K^\top\nabla v$, write
\begin{equation}\label{eq:Ac}
 S_0=S\circ\Phi,\quad q_0=q\circ\Phi,\qquad
 A=KS_0K^\top,\qquad
 c=A^\top\nabla\log J+Kq_0.
\end{equation}
Equations \eqref{eq:cancellation} and \eqref{eq:piola} yield
\begin{equation}\label{eq:transformed}
 v_t=J^{-1}\diver(JA\nabla v)+(Kq_0)\cdot\nabla v
     =\diver(A\nabla v)+c\cdot\nabla v.
\end{equation}
The transpose on $A$ in \eqref{eq:Ac} is needed when $a$ is nonsymmetric.
All factors in $c$ are bounded pathwise. Smooth diffeomorphisms preserve
spatial null sets; Fubini's theorem therefore preserves the bounds on
$S_0$ almost everywhere in spacetime. This is an integrated
$H^{-1}$ identity as well as a distributional equation. We have taken
no time derivative of $J,K$ or $b$.

\subsection{Return to the original coordinates}

\begin{proof}[Proof of Theorem~\ref{thm:main}]
On the common probability-one event above, the transformed equation
\eqref{eq:transformed} satisfies every hypothesis of Lemma~\ref{lem:frame}.
It follows that $v\in C^{\alpha/2,\alpha}([\tau,T]\times\T^d)$ with
one deterministic
\[
 \alpha=\alpha(d,\nu/4,4\Lambda_0)\in(0,1),
 \qquad \Lambda_0=M+\tfrac12M^2,
\]
where $\alpha(d,\lambda,\Lambda)$ is the choice in Lemma~\ref{lem:unit}.
The inverse-flow estimate gives, with a finite random constant,
\[
 \operatorname{dist}(\Psi_t(x),\Psi_s(y))
 \le C_\omega\bigl(\operatorname{dist}(x,y)+|t-s|^{1/4}\bigr).
\]
Since $u(t,x)=v(t,\Psi_t(x))$ in $L^2$ at every time, its continuous
representative satisfies
\begin{align*}
 |u(t,x)-u(s,y)|
 &\le C_\omega\bigl(|t-s|^{\alpha/2}
       +\operatorname{dist}(x,y)^\alpha+|t-s|^{\alpha/4}\bigr)\\
 &\le C'_\omega\bigl(|t-s|^{\alpha/4}
       +\operatorname{dist}(x,y)^{\alpha/2}\bigr)
\end{align*}
on the compact strip. Thus $\gamma=\alpha/2$ proves the assertion.
The terminal value is the actual $L^2$ trace, by the upper-trace argument
in Lemma~\ref{lem:unit} and the continuity of the pullbacks.

The representative and the membership event can be chosen measurably.
Apply a fixed sequence of periodic spatial mollifiers to the original
$L^2$-valued process. On each regular path these mollifications converge
uniformly on a compact strip to its unique continuous representative.
Values at rational spacetime points and countable limits and suprema
then determine its H\"older norm and the event that this norm is finite.
No measurable choice of the random affine charts is needed.
The same pathwise argument works on every positive-time strip; a
countable exhaustion by rational lower times gives simultaneous interior
regularity if desired.
\end{proof}

In the conclusion, $\gamma=\alpha/2$ is fixed by $d,\nu,M$, and
$C^{\gamma/2,\gamma}$ means temporal exponent $\alpha/4$ and spatial
exponent $\alpha/2$. The assumptions are precisely
\eqref{eq:bound}--\eqref{eq:smooth} and the variational interpretation of
\eqref{eq:spde}. The local radii, covering constants and H\"older norm
may depend on the realization; the exponent does not. The argument uses
neither commutation of the transport fields nor volume preservation.

\paragraph{AI assistance.}
OpenAI Codex assisted with the literature search, proof development,
manuscript preparation and separate-context automated reviews.
These automated reviews are not independent human peer review or formal
verification. Responsibility for the submitted mathematical claims rests
with the author.

\begin{thebibliography}{9}
\bibitem{ASV}
Antonio Agresti, Max Sauerbrey and Mark Veraar,
\emph{A stochastic flow approach to De Giorgi--Nash--Moser estimates for
SPDEs with smooth transport noise}, arXiv:2511.12692v1 (2025),
\url{https://arxiv.org/pdf/2511.12692v1}.
The theorem and equation locators here refer to the PDF dated November 18,
2025: distributional composition (2.33) and It\^o--Wentzell formula (2.34).
The HTML manuscript dated August 24, 2026 numbers these formulas
(2.37) and (2.38).

\bibitem{Liao}
Naian Liao, \emph{Remarks on parabolic De Giorgi classes},
Annali di Matematica Pura ed Applicata \textbf{200} (2021), 2361--2384.
\href{https://doi.org/10.1007/s10231-021-01084-8}{doi:10.1007/s10231-021-01084-8}.
\end{thebibliography}
\end{document}
